\chapter{Ingreso al sistema}

\section{Iniciar sesión}

\captura{login}{Pantalla de inicio de sesión}

\begin{enumerate}
  \item Toque su foto en la fila de arriba, o escriba su \textbf{usuario} (por ejemplo, \texttt{a.romero}).
  \item Escriba su \textbf{contraseña}. El ojo del campo la muestra u oculta.
  \item Presione \boton{Ingresar}.
\end{enumerate}

Si la contraseña es incorrecta aparece un aviso rojo y el campo se vacía para volver a escribirla. Un empleado dado de baja no puede ingresar hasta que lo reactiven.

\section{Modo DIRECTORIO}

\captura{login-directorio}{Modo DIRECTORIO}

El dueño ingresa con el botón \boton{Modo directorio}: el usuario queda como \texttt{DIRECTORIO} y solo se escribe la contraseña de administración general. Para volver al ingreso normal se usa \boton{Modo empleado}.

El DIRECTORIO cambia su contraseña con el botón de la llave, junto a su nombre arriba a la derecha. Pide la contraseña actual y la nueva dos veces.

\section{Recuperar la contraseña}

\captura[0.8]{recuperar-contrasena}{Recuperar la contraseña}

\begin{enumerate}
  \item En el inicio de sesión toque \boton{¿Olvidaste tu contraseña?}.
  \item Escriba el correo de su cuenta y presione \boton{Enviar código}.
  \item Revise su correo: llega un código de 6 dígitos que vale 10 minutos. Escríbalo y presione \boton{Verificar}.
  \item Escriba la nueva contraseña dos veces (mínimo 8 caracteres, con letras y números) y presione \boton{Guardar contraseña}.
\end{enumerate}

\begin{importante}
El código admite 5 intentos. Si se equivoca más veces, pida un código nuevo. Si su cuenta está dada de baja, el sistema no envía el código: pida a administración que lo reactive.
\end{importante}

\section{Cerrar sesión}

Presione \boton{Cerrar sesión}, arriba a la derecha (en el celular está al final del menú). Hágalo siempre al terminar el turno, para que nadie registre pedidos a su nombre.
